% ---- score2 standard preamble (from: python3 score2_tools.py preamble) ----
\documentclass[10pt]{extarticle}
\usepackage[a4paper,margin=1.5cm]{geometry}
\usepackage[T1]{fontenc}
\usepackage{amsmath,amssymb}
\usepackage{graphicx}
\usepackage[export]{adjustbox}
\usepackage{multicol}
\usepackage{enumitem}
\usepackage{booktabs,array,tabularx}
\setlength{\columnsep}{0.7cm}
\setlength{\parindent}{0pt}
\setlength{\parskip}{3pt}
\pagestyle{empty}
\makeatletter
% figure registry (figdefs.tex is generated by the extract step)
\newcommand{\qfigdef}[4]{\@namedef{qf@file@#1}{#2}\@namedef{qf@w@#1}{#3}\@namedef{qf@wide@#1}{#4}}
\newcommand{\qf@check}[1]{\@ifundefined{qf@file@#1}{\PackageError{score2}{Unknown figure id #1}{}}{}}
\newcommand{\qf@img}[4][]{\includegraphics[width=#2\dimexpr\@nameuse{qf@w@#4}pt\relax,max width=#3,max height=0.45\textheight,#1]{\@nameuse{qf@file@#4}}}
% Full-width questions: finalize wraps every question that contains a wide figure in \qfullbegin...\qfullend,
% so the WHOLE question (stem, figures, options) is set across the page in reading order, then the two columns
% resume. The question's head - its number, stem and first figure (\qfullsplit marks the end of the head) -
% is kept on one page, so a number is never left at a page foot with its figure overleaf; the rest may break
% across pages normally, so no big gaps are left. Without a \qfullsplit, the whole question is kept together
% if it is under 45% of a page. Never write these yourself.
\newif\ifqf@full
\newif\ifqf@boxed
\newsavebox{\qf@box}
\newcommand{\qf@open}{\global\qf@fulltrue\par\global\qf@boxedtrue
  \setbox\qf@box\vbox\bgroup\hsize\textwidth\linewidth\textwidth}
\newcommand{\qf@emit}[1]{\par\egroup\global\qf@boxedfalse
  \ifdim\dimexpr\ht\qf@box+\dp\qf@box\relax<#1\textheight\noindent\usebox{\qf@box}\par\else\unvbox\qf@box\fi}
\newcommand{\qfullsplit}{\ifqf@boxed\qf@emit{0.85}\fi}
\newcommand{\qf@close}{\ifqf@boxed\qf@emit{0.45}\fi\global\qf@fullfalse}
\newcommand{\qfullbegin}{\end{multicols}\qf@open}
\newcommand{\qfullend}{\qf@close\begin{multicols}{2}}
\newcommand{\qfullstart}{\qf@open}
\newcommand{\qfulldone}{\qf@close}
% \qfig[scale]{id}: block figure, centred, directly where it is written (size set by the tool).
% Fallback if a wide figure is not inside a full-width question: break out of the columns in place - no floats.
\newcommand{\qfig}[2][1]{\qf@check{#2}%
  \ifqf@full\qf@blk{#1}{#2}\else
  \ifnum\@nameuse{qf@wide@#2}=1 \expandafter\@firstoftwo\else\expandafter\@secondoftwo\fi
  {\end{multicols}\begin{center}\qf@img{#1}{\textwidth}{#2}\end{center}\begin{multicols}{2}}%
  {\qf@blk{#1}{#2}}\fi}
\newcommand{\qf@blk}[2]{\par\smallskip{\centering\qf@img{#1}{\linewidth}{#2}\par}\smallskip}
% \qfigi[scale]{id}: inline figure (inside option lists / tables / minipages); never breaks columns.
\newcommand{\qfigi}[2][1]{\qf@check{#2}\qf@img[valign=c]{#1}{\linewidth}{#2}}
\makeatother
% structure macros - lint reads these to check completeness and figure placement
\newcommand{\papersec}[1]{\par\bigskip{\centering\Large\bfseries PAPER-#1\par}\medskip}
\newcommand{\partsec}[2]{\par\medskip\noindent{\large\bfseries PART-#1: #2}\par\smallskip}
\newcommand{\sect}[1]{\par\smallskip\noindent\textbf{#1}\par}
\newcommand{\qn}[1]{\par\smallskip\noindent\textbf{#1.}~}
\newcommand{\paraq}[2]{\par\smallskip\noindent\textbf{Paragraph for Questions #1 and #2}\par}
\newenvironment{qoptions}{\begin{itemize}[leftmargin=1.8em,itemsep=1pt,topsep=2pt,label={}]}{\end{itemize}}
\newcommand{\opt}[2]{\item[(#1)] #2}
% figure registry (generated by score2_tools.py finalize)
\qfigdef{p1_phy_s1_q04_1}{TOAS - Test 02 (13 Feb 2022)-img/p1_phy_s1_q04_1.png}{75.1}{0}
\qfigdef{p1_phy_s1_q05_1}{TOAS - Test 02 (13 Feb 2022)-img/p1_phy_s1_q05_1.png}{115.4}{0}
\qfigdef{p1_phy_s2_q07_1}{TOAS - Test 02 (13 Feb 2022)-img/p1_phy_s2_q07_1.png}{130.6}{0}
\qfigdef{p1_phy_s2_q08_1}{TOAS - Test 02 (13 Feb 2022)-img/p1_phy_s2_q08_1.png}{129.4}{0}
\qfigdef{p1_phy_s2_q09_1}{TOAS - Test 02 (13 Feb 2022)-img/p1_phy_s2_q09_1.png}{157.0}{0}
\qfigdef{p1_phy_s2_q11_1}{TOAS - Test 02 (13 Feb 2022)-img/p1_phy_s2_q11_1.png}{207.5}{0}
\qfigdef{p1_phy_s3_q04_1}{TOAS - Test 02 (13 Feb 2022)-img/p1_phy_s3_q04_1.png}{167.5}{0}
\qfigdef{p1_phy_s3_q05_1}{TOAS - Test 02 (13 Feb 2022)-img/p1_phy_s3_q05_1.png}{245.2}{0}
\qfigdef{p1_phy_s3_q06_1}{TOAS - Test 02 (13 Feb 2022)-img/p1_phy_s3_q06_1.png}{126.3}{0}
\qfigdef{p1_chem_s1_q03_A}{TOAS - Test 02 (13 Feb 2022)-img/p1_chem_s1_q03_A.png}{45.6}{0}
\qfigdef{p1_chem_s1_q03_B}{TOAS - Test 02 (13 Feb 2022)-img/p1_chem_s1_q03_B.png}{44.8}{0}
\qfigdef{p1_chem_s1_q03_C}{TOAS - Test 02 (13 Feb 2022)-img/p1_chem_s1_q03_C.png}{33.5}{0}
\qfigdef{p1_chem_s1_q03_D}{TOAS - Test 02 (13 Feb 2022)-img/p1_chem_s1_q03_D.png}{70.1}{0}
\qfigdef{p1_chem_s1_q05_A}{TOAS - Test 02 (13 Feb 2022)-img/p1_chem_s1_q05_A.png}{128.6}{0}
\qfigdef{p1_chem_s1_q05_B}{TOAS - Test 02 (13 Feb 2022)-img/p1_chem_s1_q05_B.png}{132.5}{0}
\qfigdef{p1_chem_s1_q05_C}{TOAS - Test 02 (13 Feb 2022)-img/p1_chem_s1_q05_C.png}{131.8}{0}
\qfigdef{p1_chem_s1_q05_D}{TOAS - Test 02 (13 Feb 2022)-img/p1_chem_s1_q05_D.png}{144.1}{0}
\qfigdef{p1_chem_s2_q09_A}{TOAS - Test 02 (13 Feb 2022)-img/p1_chem_s2_q09_A.png}{235.3}{0}
\qfigdef{p1_chem_s2_q09_B}{TOAS - Test 02 (13 Feb 2022)-img/p1_chem_s2_q09_B.png}{244.2}{0}
\qfigdef{p1_chem_s2_q09_C}{TOAS - Test 02 (13 Feb 2022)-img/p1_chem_s2_q09_C.png}{168.9}{0}
\qfigdef{p1_chem_s2_q09_D}{TOAS - Test 02 (13 Feb 2022)-img/p1_chem_s2_q09_D.png}{245.2}{0}
\qfigdef{p1_chem_s2_q10_1}{TOAS - Test 02 (13 Feb 2022)-img/p1_chem_s2_q10_1.png}{153.7}{0}
\qfigdef{p1_chem_s2_q11_1}{TOAS - Test 02 (13 Feb 2022)-img/p1_chem_s2_q11_1.png}{198.1}{0}
\qfigdef{p1_chem_s3_q02_1}{TOAS - Test 02 (13 Feb 2022)-img/p1_chem_s3_q02_1.png}{113.7}{0}
\qfigdef{p1_chem_s3_q03_1}{TOAS - Test 02 (13 Feb 2022)-img/p1_chem_s3_q03_1.png}{185.2}{0}
\qfigdef{p1_chem_s3_q04_1}{TOAS - Test 02 (13 Feb 2022)-img/p1_chem_s3_q04_1.png}{260.6}{1}
\qfigdef{p1_math_s2_q12_1}{TOAS - Test 02 (13 Feb 2022)-img/p1_math_s2_q12_1.png}{198.5}{0}
\qfigdef{p1_math_s2_q12_2}{TOAS - Test 02 (13 Feb 2022)-img/p1_math_s2_q12_2.png}{103.1}{0}
\qfigdef{p1_math_s2_q12_3}{TOAS - Test 02 (13 Feb 2022)-img/p1_math_s2_q12_3.png}{132.3}{0}
\qfigdef{p1_math_s2_q12_4}{TOAS - Test 02 (13 Feb 2022)-img/p1_math_s2_q12_4.png}{139.9}{0}
\qfigdef{p1_math_s2_q12_5}{TOAS - Test 02 (13 Feb 2022)-img/p1_math_s2_q12_5.png}{183.0}{0}
\raggedcolumns
\makeatletter\@beginparpenalty=10000\makeatother
\newcolumntype{L}[1]{>{\raggedright\arraybackslash}p{#1}}
\begin{document}
\begin{center}{\Large\bfseries ALLEN TOAS --- Test 02 (13 Feb 2022)}\end{center}
\begin{multicols}{2}
\partsec{1}{PHYSICS}
\sect{SECTION-I(i)}
\qn{1} A particle is released from rest from a certain height $H=80$ m. Due to the wind the particle gathers the horizontal velocity component $v_x=ay$ where $a=3/8\ \mathrm{s^{-1}}$ and $y$ is the vertical displacement of the particle from the point of release. Find the horizontal displacement of the particle when it strikes the ground. (Take : $g=10\ \mathrm{m/s^2}$)
\begin{qoptions}
\opt{A}{50 m}
\opt{B}{40 m}
\opt{C}{30 m}
\opt{D}{None of these}
\end{qoptions}
\qn{2} The displacement of particle executing SHM is given by $x=A\sin(\omega t+\theta)$, where $\theta$-initial phase of motion, $A$-Amplitude of motion and $T$-Time period. The magnitude of average velocity during motion of particle between $t=0$ to $t=T/4$ is :
\begin{qoptions}
\opt{A}{$2\dfrac{A\omega}{\pi}$}
\opt{B}{$2\dfrac{A\omega}{\pi}|\cos\theta|$}
\opt{C}{$2\dfrac{A\omega}{\pi}|\cos\theta-\sin\theta|$}
\opt{D}{$2\dfrac{A\omega}{\pi}|\cos\theta+\sin\theta|$}
\end{qoptions}
\qn{3} Two particles $A$ and $B$ are performing SHM along $x$ and $y$-axis respectively with equal amplitude and frequency of 2 cm and 1 Hz respectively. Equilibrium positions of the particles $A$ and $B$ are at the co-ordinates (3, 0) cm and (0, 4) cm respectively. At $t=0$, $B$ is at its equilibrium position and moving towards the origin, while $A$ is nearest to the origin and moving away from the origin. If smallest separation between $A$ and $B$ is $d_{\min}$ and maximum separation is $d_{\max}$ then $d_{\max}-d_{\min}=$ \underline{\hspace{2.5em}} ?
\begin{qoptions}
\opt{A}{4 cm}
\opt{B}{5 cm}
\opt{C}{2 cm}
\opt{D}{1.5 cm}
\end{qoptions}
\qn{4} In the figure shown, for an angle of incidence $i$ at the top surface, what is the minimum refractive index needed for total internal reflection at the vertical face?
\qfig{p1_phy_s1_q04_1}
\begin{qoptions}
\opt{A}{$\sqrt{1-\sin^2i}$}
\opt{B}{$\sqrt{1+\sin^2i}$}
\opt{C}{$\sqrt{1-\cos^2i}$}
\opt{D}{$\sqrt{1+\cos^2i}$}
\end{qoptions}
\qn{5} A point $A$ is kept on the axis of a converging lens of focal length 10 cm as shown in the figure. A plane mirror $M$ is placed inclined at $45^\circ$ to the lens axis at a distance of 10 cm to the right of the lens. The co-ordinates of the image of $A$ formed by this combination are (take the centre of the lens as the origin and the lens axis as the axis of $x$). (Consider only two optical phenomenon)
\qfig{p1_phy_s1_q05_1}
\begin{qoptions}
\opt{A}{($-10$ cm, 20 cm)}
\opt{B}{($-10$ cm, $-20$ cm)}
\opt{C}{(10 cm, $-20$ cm)}
\opt{D}{(10 cm, 20 cm)}
\end{qoptions}
\qn{6} A projectile is required to hit a target whose co-ordinates relative to horizontal and vertical axes through the point of projection are $(\alpha,\beta)$. If the gun velocity is $\sqrt{2g\alpha}$, then the condition when it is impossible to hit the target, is
\begin{qoptions}
\opt{A}{$3\alpha>4\beta$}
\opt{B}{$4\beta>3\alpha$}
\opt{C}{$3\beta>2\alpha$}
\opt{D}{$4\alpha>3\beta$}
\end{qoptions}
\sect{SECTION-I(ii)}
\qn{7} A parallel beam of light is incident normally on the flat surface of a hemisphere of radius 6cm and refractive index 1.5, placed in air as shown in figure (i). Assume paraxial ray approximation. Choose the correct option(s) :-
\qfig{p1_phy_s2_q07_1}
\begin{qoptions}
\opt{A}{The rays are focussed at 12 cm from the point $P$ to the right, in the situation as shown in figure (i)}
\opt{B}{The rays are focussed at 16 cm from the point $P$ to the right, in the situation as shown in figure (i)}
\opt{C}{If the rays are incident at the curved surface (figure (ii)) then these are focused at distance 18 cm from point $P$ to the right}
\opt{D}{If the rays are incident at the curved surface (figure (ii)) then these are focused at distance 14 cm from point $P$ to the right}
\end{qoptions}
\qn{8} Two points $P$ and $Q$ lie on either side of an axis $XY$ as shown. It is desired to produce an image at $Q$ for a real object placed at $P$ using a spherical mirror, with $XY$ as the optic axis. The mirror must be
\qfig{p1_phy_s2_q08_1}
\begin{qoptions}
\opt{A}{converging and positioned to the left of $P$}
\opt{B}{diverging and positioned to the left of $P$}
\opt{C}{converging and positioned to the right of $Q$}
\opt{D}{diverging and positioned to the right of $Q$}
\end{qoptions}
\qn{9} Figure shows block 1 of mass 0.5 kg sliding to the right over a frictionless elevated surface at a speed of 12 m/s. The block undergoes an elastic collision with stationary block 2, which is attached to a spring of spring constant 100 N/m (Assume that spring doesn't affect collision). After the collision, block 2 oscillates in SHM with a time period of 0.628 s and block 1 slides to the opposite end of the elevated surface landing at a distance $d$ from the base of that surface after falling through a height $h=4.90$ m. Then which of the following is (are) \textbf{CORRECT}: (take $g=9.8\ \mathrm{m/s^2}$ and $\pi=3.14$)
\qfig{p1_phy_s2_q09_1}
\begin{qoptions}
\opt{A}{The value of $d=2$ m}
\opt{B}{Spring potential energy is $<35$ J}
\opt{C}{Maximum kinetic energy of block 2 is 32 J}
\opt{D}{Speed of block 1 when it hits the ground is 4 m/s}
\end{qoptions}
\qn{10} Two boats are at sea and moving continuously in straight lines. At 11 AM boat 1 was 40 km North of a small island and was moving at 15 km/h towards east. The boat 2 at 7 AM of the same day was 100 km away east of the same island and moved at a speed of 15 km/h towards the South. Which of the following is (are) \textbf{CORRECT}?
\begin{qoptions}
\opt{A}{The minimum separation between the boats is $100\sqrt2$ km}
\opt{B}{The minimum separation between the boats is less than $100\sqrt2$ km}
\opt{C}{The boats will be found at minimum separation at 9 AM}
\opt{D}{The boats will be found at minimum separation at 11 AM}
\end{qoptions}
\qn{11} In the figure shown, a ball is launched with a velocity of magnitude 10 m/s at an angle $53^\circ$ with the horizontal. The launch point is at the base of a ramp of horizontal length $d_1$ and height $d_2$. A plateau (flat surface) is located at the top of the ramp. Which of the following option(s) is(are) \textbf{CORRECT}? (Take : $g=10\ \mathrm{m/s^2}$)
\qfig{p1_phy_s2_q11_1}
\begin{qoptions}
\opt{A}{If $d_1=12$ m and $d_2=5$m, the ball lands of ramp.}
\opt{B}{If $d_1=6$m and $d_2=5$m, ball lands on plateau.}
\opt{C}{If $d_1=6$m and $d_2=3$m, the displacement of the ball from the point of projection to the point where it hits the surface is $3\sqrt5$m}
\opt{D}{If $d_1=6$m and $d_2=1$m, displacement of the particle is $>3\sqrt5$m}
\end{qoptions}
\qn{12} A particle moves along the x-axis according to the law $x=a\sin^2\left(\omega t-\dfrac\pi4\right)$
\begin{qoptions}
\opt{A}{The amplitude of the oscillation is $a$}
\opt{B}{The time period of oscillations is $\dfrac\pi\omega$}
\opt{C}{The velocity ($v$) of the particle is given by $v^2=4\omega^2x(a-x)$}
\opt{D}{The velocity ($v$) of the particle is given by $v^2=2\omega^2x(a-x)$}
\end{qoptions}
\sect{SECTION-II}
\qn{1} When the plane surface of plano-convex lens is silvered, the lens behaves like a concave mirror of focal length 30 cm, but when the convex surface is silvered, it behaves like a concave mirror of focal length 10 cm. The refractive index of material of lens is $\dfrac{P}{2}$. What is the value of P?
\qn{2} A bird flies with a uniform speed of 10 km/h and a car moves with a uniform speed of 8 km /h. Both start from $B$ towards $A$ ($BA=40$ km) at the same instant. The bird having reached $A$, flies back to meet the approaching car. As soon as it reaches the car, it flies back to $A$. Assume that bird immediately gains the same speed while turning back. The distance of the car from point $A$ is $\dfrac{16}{\alpha}$ when the fly reaches point $A$ for the 3rd time. The value of $\alpha$ is :
\qn{3} A ray of light passes through a prism in principle plane. The angle of prism is $\alpha$. Angle of deviation and angle of incidence both are equal to $2\alpha$. If $\mu$ is the refractive index of the prism then it can be found that $\cos^2\alpha=\dfrac{\mu^n-1}{2x}$ where $n$ is a positive integer. The value of $x$ is:
\qn{4} Figure shows an overhead view of a corridor with a plane mirror $M$ mounted at one end. A burglar $B$ sneaks along the corridor directly towards the centre of the mirror. If $d=4$m, how far (in meter) from the mirror is the burglar when he will be first seen in the mirror by the security guard $S$ ?
\qfig{p1_phy_s3_q04_1}
\qn{5} A metal plate is mounted on an axle through its centre of mass. A spring with spring constant $k=2000$ N/m connects a wall with a point on the rim a distance $r=2.5$ cm from the centre of mass. Initially the spring is at its rest length. If the plate is rotated by $7^\circ$ and released, it rotates about the axle in SHM with its angular position given by the graph shown. The horizontal axis scale is set by $t_s=31.4$ ms. If the moment of inertia of the plate about the centre of mass is $x\times10^{-6}\ \mathrm{kg\,m^2}$. Then the value of $x$ is:
\qfig{p1_phy_s3_q05_1}
\qn{6} A diverging lens of focal length 10 cm is placed 10 cm in front of a plane mirror as shown in the figure. Light from a very far away source falls on the lens. The final image is at a distance "$V$" cm from the mirror. Write the value of $V$ (positive) :
\qfig{p1_phy_s3_q06_1}
\partsec{2}{CHEMISTRY}
\sect{SECTION-I(i)}
\qn{1} On mixing 15.0 ml of ethyl alcohol of density 0.792 g $\mathrm{ml^{-1}}$ with 15 ml of pure water at $4^\circ$C, the resulting solution is found to have a density of 0.924 g $\mathrm{ml^{-1}}$. The percentage contraction in volume is :
\begin{qoptions}
\opt{A}{8 \%}
\opt{B}{2 \%}
\opt{C}{3 \%}
\opt{D}{4 \%}
\end{qoptions}
\qn{2} Which of the following is intrinsic semiconductor?
\begin{qoptions}
\opt{A}{Si}
\opt{B}{Si doped with P}
\opt{C}{GaAs}
\opt{D}{$\mathrm{Fe_{0.93}O_{1.00}}$}
\end{qoptions}
\qn{3} Which of the following compound is optically active with $\mathrm{C_2}$-axis of symmetry ?
\begin{qoptions}
\opt{A}{\qfigi{p1_chem_s1_q03_A}}
\opt{B}{\qfigi{p1_chem_s1_q03_B}}
\opt{C}{\qfigi{p1_chem_s1_q03_C}}
\opt{D}{\qfigi{p1_chem_s1_q03_D}}
\end{qoptions}
\qn{4} Select correct statement(s) among the following
\begin{qoptions}
\opt{A}{Polar protic solvent favours $\mathrm{S_N2}$ mechanism more than polar aprotic solvent generally}
\opt{B}{Nucleophilic order of halides in polar protic solvent is\\ $\mathrm{I^{\ominus}>Br^{\ominus}>Cl^{\ominus}>F^{\ominus}}$}
\opt{C}{Always perfect racemisation take place in $\mathrm{S_N1}$ mechanism}
\opt{D}{Inversion of configuration in $\mathrm{S_N2}$ mechanism always produces R to S configuration or S to R configuration}
\end{qoptions}
\qn{5} Which of the following graph correctly depicts the property mentioned against the indicated species.
\begin{qoptions}
\opt{A}{\qfigi{p1_chem_s1_q05_A}}
\opt{B}{\qfigi{p1_chem_s1_q05_B}}
\opt{C}{\qfigi{p1_chem_s1_q05_C}}
\opt{D}{\qfigi{p1_chem_s1_q05_D}}
\end{qoptions}
\qn{6} The specie that show maximum number of stereo Isomers among the following.
\begin{qoptions}
\opt{A}{$\mathrm{[Ni(H_2O)_2(en)_2]^{2+}}$}
\opt{B}{$\mathrm{[Co(NH_3)_3(NO_2)_3]}$}
\opt{C}{$\mathrm{[Pt(py)_2(NH_3)_2Cl_2]^{2+}}$}
\opt{D}{$\mathrm{[CoBr_2Cl_2(NO_2)(SCN)]^{3-}}$}
\end{qoptions}
\sect{SECTION-I(ii)}
\qn{7} A metal `M' crystallizes in two crystal systems fcc and hcp. If $d$-density, $z$ - number of atoms per unit cell and $a,b,c$ are edge lengths of the unit cell along the crystallographic axes then. Select the correct options :
\begin{qoptions}
\opt{A}{$\dfrac{d_{\mathrm{fcc}}}{d_{\mathrm{hcp}}}\ne1$}
\opt{B}{$\dfrac{a_{\mathrm{fcc}}}{c_{\mathrm{hcp}}}=\dfrac{\sqrt3}{2}$}
\opt{C}{$\dfrac{d_{\mathrm{fcc}}}{d_{\mathrm{hcp}}}=1$}
\opt{D}{$\dfrac{z_{\mathrm{fcc}}}{z_{\mathrm{hcp}}}=\dfrac23$}
\end{qoptions}
\qn{8} The correct information(s) about $\mathrm{CaF_2}$ crystal (Fluorite structure) is/are.
\begin{qoptions}
\opt{A}{C.N. of $\mathrm{Ca^{2+}}$ ion is 8}
\opt{B}{C.N. of $\mathrm{F^-}$ ion is 8}
\opt{C}{$\mathrm{F^-}$ ions occupy all the tetrahedral voids formed by CCP arrangement of $\mathrm{Ca^{2+}}$ ions}
\opt{D}{$\mathrm{Ca^{2+}}$ ions occupy all the cubic voids formed by simple cubic arrangement of $\mathrm{F^-}$ ions}
\end{qoptions}
\qn{9} Which of the following reaction is showing correct major product :
\begin{qoptions}
\opt{A}{\qfigi{p1_chem_s2_q09_A}}
\opt{B}{\qfigi{p1_chem_s2_q09_B}}
\opt{C}{\qfigi{p1_chem_s2_q09_C}}
\opt{D}{\qfigi{p1_chem_s2_q09_D}}
\end{qoptions}
\qn{10} \qfigi{p1_chem_s2_q10_1}\par
Selelct correct statement form the following about major product A, B and C :
\begin{qoptions}
\opt{A}{A is $3^{\mathrm o}$ alcohol and show optical isomersim}
\opt{B}{B is $1^{\mathrm o}$ alcohol and does not show optical isomersim}
\opt{C}{C is $2^{\mathrm o}$ alcohol and shows optical isomersim}
\opt{D}{A and B are positional isomer of each other}
\end{qoptions}
\qn{11} One of the isomer of $[\mathrm{Pd(NH_2{-}CH(CH_3){-}CO_2^-)_2}]$ is given as
\qfig{p1_chem_s2_q11_1}
The correct statement(s) regarding the above compound is/are.
\begin{qoptions}
\opt{A}{Total stereoisomers will be 6}
\opt{B}{The ligand $\mathrm{(NH_2{-}CH(CH_3){-}COO^-)}$ is Ambidentate ligand}
\opt{C}{The molecule is planar since $\mathrm{Pd^{2+}}$ shows square planar geometry around it}
\opt{D}{Exactly two Enantiomeric pairs will be formed}
\end{qoptions}
\qn{12} The compound $\mathrm{AB_3}$ has more dipole moment than $\mathrm{AC_3}$. The compound $\mathrm{AC_3}$ on action of superheated steam under high pressure yield an acidic oxide via $\mathrm{S_N1}$ mechanism. $\mathrm{AB_3}$ has tendency to unsymmetrically cleave bridge bond of $\mathrm{B_2H_6}$. Which is/are correct in view of above information.
\begin{qoptions}
\opt{A}{The compound ``BC'' can be produced by hydrolysis of $\mathrm{C_2}$.}
\opt{B}{The compound $\mathrm{BA_3}$ has less B--A bond length compared to C--A bond length in $\mathrm{CA_3}$.}
\opt{C}{$\mathrm{AC_3}$ show back bonding.}
\opt{D}{A can form neutral oxide}
\end{qoptions}
\sect{SECTION-II}
\qn{1} A compound has cubical unit cell in which `A' atoms are found at six corners and `B' atoms are present at remaining corners and at face centres which are not opposite to each other and `C' atoms are present at body centres and remaining face centres. Calculate density of the compound in $\mathrm{amu/\AA^3}$.\\
Given : A = 20 amu, B = 60 amu, C = 80 amu and edge length = 2\AA.
\qn{2} A metal chloride of empirical formula MCl is formed (where M is metal). In an experiment weight of metal chloride has following variation according to weight of chlorine.
\qfig{p1_chem_s3_q02_1}
Then atomic weight of metal (M) will be :
\qn{3} How many of the following compound will form immediate turbidity with ($\mathrm{HCl/ZnCl_2}$):
\qfig{p1_chem_s3_q03_1}
\qfullbegin
\qn{4} Number of molecules show optical isomerism :
\qfig{p1_chem_s3_q04_1}
\qfullend
\qn{5} Which of the following compounds/species are incorrectly matched against the indicated names.\\
(a) $\mathrm{[PdI_2(ONO)_2(H_2O)_2]}\to$ Diaquadiiodidodinitrito-Npalladium(IV)\\
(b) $\mathrm{Fe(\pi\text{-}C_5H_5)_2}\to$ Bis($\eta^1$--cyclopentadienyl)iron(II)\\
(c) $\mathrm{H_2[PtCl_6]}\to$ Hexachloridoplatinate(IV) hydride\\
(d) $\mathrm{[Al(OH)(H_2O)_5]SO_4}\to$ Pentaaquahydroxidoaluminium(III) sulphate\\
(e) $\mathrm{K_2[OsCl_5N]}\to$ Potassium nitridopentachloridoosmate(VI)\\
(f) $\mathrm{Cr(C_6H_6)_2}\to$ Bis($\eta$6--benzene)chromium(0)
\qn{6} On the basis of the following observations made with aqueous solutions of given complexes, assigning secondary valency of each metal, write the sum of secondary valencies of all given complexes as your answer.
\par\nobreak
{\small\setlength{\tabcolsep}{3pt}\renewcommand{\arraystretch}{1.25}\begin{center}
\begin{tabular}{|L{2.9cm}|L{4.6cm}|}
\hline
Formula & Moles of AgCl precipitated per mole of compound with excess $\mathrm{AgNO_3}$ solution\\
\hline
(i)\ \ $\mathrm{PdCl_2.4NH_3}$ & 2 moles\\
\hline
(ii)\ \ $\mathrm{NiCl_2.6NH_3}$ & 2 moles\\
\hline
(iii)\ \ $\mathrm{CoCl_3.4NH_3}$ & 1 mole\\
\hline
(iv)\ \ $\mathrm{PtCl_2.2NH_3}$ & Zero mole\\
\hline
\end{tabular}\end{center}}
\partsec{3}{MATHEMATICS}
\sect{SECTION-I(i)}
\qn{1} Consider a hyperbola $x^2-y^2-2x+4y-7=0$. If its pair of asymptotes meet two tangents each of slope 2 at $P,Q$ and $R,S$ (in order) then choose the INCORRECT statement :
\begin{qoptions}
\opt{A}{$PQ=\dfrac{4\sqrt5}{\sqrt3}$}
\opt{B}{Latus rectum of hyperbola is 4}
\opt{C}{area of quadrilateral $PQRS$ is 12 sq.units}
\opt{D}{area of quadrilateral $PQRS$ is 16 sq.units}
\end{qoptions}
\qn{2} Consider $f:R\to R$, $f(x)=\dfrac{3x^2+ax+b}{x^2+1}$ where $a,b\in R$. If the range of this function is $[-4,3)$ then $a+b$ is equal to
\begin{qoptions}
\opt{A}{0}
\opt{B}{$-4$}
\opt{C}{$-1$}
\opt{D}{3}
\end{qoptions}
\qn{3} Given that $f(x)$ is a function defined on $\mathbf R$, satisfying $f(1)=1$ and for any $x\in\mathbf R$, $f(x+5)\ge f(x)+5$ and $f(x+1)\le f(x)+1$. If $g(x)=f(x)+5-x$, then $g(2021)$ is equal to -
\begin{qoptions}
\opt{A}{3}
\opt{B}{4}
\opt{C}{5}
\opt{D}{6}
\end{qoptions}
\qn{4} Let $f:R\to R$ and $f(x)$ be a polynomial with real coefficient such that $(x^2+x+1)\,f(x-1)=(x^2-x+1)\,f(x)\ \forall\ x\in R$ and $f(0)=1$, then which of the following is false -
\begin{qoptions}
\opt{A}{$f(x)$ is many one function}
\opt{B}{$f(x)$ has no real roots}
\opt{C}{range of $f(x)$ is $\left[\dfrac12,\infty\right)$}
\opt{D}{$f(x)$ is neither even nor odd}
\end{qoptions}
\qn{5} If $[x]^3+[y]^3+[z]^3-3[x][y][z]=4$; where $x,y,z>0$ (where [.] denotes GIF); then minimum possible value of $(x+y+z)$ is
\begin{qoptions}
\opt{A}{2}
\opt{B}{3}
\opt{C}{4}
\opt{D}{6}
\end{qoptions}
\qn{6} Solution set of equation is $\left|1-\log_{\frac16}x\right|+|\log_2x|+2=\left|3-\log_{\frac16}x+\log_{\frac12}x\right|$, $a,b\in N$, then the value of $(a+b)$ is
\begin{qoptions}
\opt{A}{5}
\opt{B}{6}
\opt{C}{7}
\opt{D}{8}
\end{qoptions}
\sect{SECTION-I(ii)}
\qn{7} Tangent is drawn at any point $(x_1,y_1)$ other than the vertex, on parabola $y^2=4ax$. If tangents are drawn from any point on this tangent to the circle $x^2+y^2=a^2$ such that all the chords of contact pass through a fixed point $(x_2,y_2)$ then
\begin{qoptions}
\opt{A}{$x_1,a,x_2$ are in GP}
\opt{B}{$\dfrac{y_1}{2},a,y_2$ are in GP}
\opt{C}{$-4,\dfrac{y_1}{y_2},\dfrac{x_1}{x_2}$ are in GP}
\opt{D}{$x_1x_2+y_1y_2=a^2$}
\end{qoptions}
\qn{8} A straight line $y=m(x-a)$, $m\ne0$ meets a parabola $y^2=4ax$ at two points $A$ and $B$. Then which of the following option(s) is/are correct
\begin{qoptions}
\opt{A}{Length $AB=4a$ when $m=1$}
\opt{B}{Length $AB=8a$ when $m=1$}
\opt{C}{Angle between normals drawn at $A$ and $B$ is $45^\circ$ when $m=1$}
\opt{D}{Angle between normals drawn at $A$ and $B$ is $90^\circ$ when $m=2$}
\end{qoptions}
\qn{9} If tangents $PQ$ and $PR$ are drawn from a point on the circle $x^2+y^2=25$ to the ellipse $\dfrac{x^2}{16}+\dfrac{y^2}{b^2}=1$, ($b<4$), so that the fourth vertex $S$ of the parallelogram $PQSR$ lies on the circumcircle of triangle $PQR$, then which of the following is/are true
\begin{qoptions}
\opt{A}{$e=\dfrac{\sqrt7}{4}$}
\opt{B}{$e=\dfrac{\sqrt5}{4}$}
\opt{C}{$b=3$}
\opt{D}{$b=\sqrt{11}$}
\end{qoptions}
\qn{10} Let ordered pair $(x,y)$ satisfy the equation $49^{x^2+y}+49^{y^2+x}=\dfrac{2}{\sqrt7}$ then which of the following is/are correct ?
\begin{qoptions}
\opt{A}{Only 1 ordered pair of $(x,y)$ is possible}
\opt{B}{More than 1 ordered pair of $(x,y)$ are possible}
\opt{C}{In one of the ordered pair $|x|+|y|=1$}
\opt{D}{In one of the ordered pair $x-y=0$}
\end{qoptions}
\qn{11} The values of $x$ satisfying the equation $x^9+\dfrac98x^6+\dfrac{27}{64}x^3-x+\dfrac{219}{512}=0$ is -
\begin{qoptions}
\opt{A}{$\dfrac{-1+\sqrt{13}}{4}$}
\opt{B}{$-\dfrac12$}
\opt{C}{$\dfrac{-1-\sqrt{13}}{4}$}
\opt{D}{$\dfrac12$}
\end{qoptions}
\qn{12} The graph of the function $y=f(x)$ is shown in the figure. Then which one of the following graph(s) of corresponding curve is/are correct ?
\qfig{p1_math_s2_q12_1}
\begin{qoptions}
\opt{A}{$|y|=\mathrm{sgn}(f(x))$\\ \qfigi{p1_math_s2_q12_2}}
\opt{B}{$|y|=\mathrm{sgn}(-f(|x|))$\\ \qfigi{p1_math_s2_q12_3}}
\opt{C}{$|y|=|f(|x|)|$\\ \qfigi{p1_math_s2_q12_4}}
\opt{D}{$y=x\,\mathrm{sgn}(f(x))$\\ \qfigi{p1_math_s2_q12_5}}
\end{qoptions}
\sect{SECTION-II}
\qn{1} Let $A=\{1,2,3,4\}$, $B=\{1,2,3,4\}$. The total number of functions $f:A\to B$ such that $f(x)\ne x\ \forall\ x\in A$ and range of $f(x)$ contains exactly three elements, is
\qn{2} Let $T_n$ be the $n^{\mathrm{th}}$ term of a series for $n\ge1$ and $T_1=1$, $T_{n+1}=2T_n+4^n$, $n\ge1$. If $S_n=$ Sum of first $n$ terms of the series and $S_n=\dfrac{a.4^n+b.2^n+1}{3}$ where $a,b\in I$, then the value of $\left|\dfrac ba\right|$ is
\qn{3} Let $a,b$ and $c$ are distinct positive numbers, none of them is equal to unity such that $\log_ba.\log_ca+\log_ab.\log_cb+\log_ac.\log_bc-\log_ba\sqrt a\log_{\sqrt c}b^{1/3}\log_ac^3=0$, then the value of $abc$ is -
\qn{4} If the fundamental period of the function $f(x)=\sin x+\tan\dfrac x2+\sin\dfrac x{2^2}+\tan\dfrac x{2^3}+\dots+\sin\dfrac x{2^7}+\tan\dfrac x{2^8}$ is $k\pi$, then $k$ is equal to -
\qn{5} Two perpendicular normals are drawn at the point $P$ and $Q$ of the ellipse $\dfrac{x^2}{16}+\dfrac{y^2}{9}=1$, intersecting each other on the major axis at a point $A$ (other than origin). If the area of $\Delta APQ$ is $\lambda$, then the value of $\lambda$ equals
\qn{6} Area enclosed within the curve $[y]=2^{[x]}$, $x\in[0,9)$ (where [$\bullet$] denotes the greatest integer function) is equal to
\end{multicols}
\end{document}
